\documentclass[t, handout, xcolor=table, aspectratio=169]{beamer}
\usetheme{CambridgeUS}

\usepackage{listings}
\usepackage{xcolor}
\usepackage{multicol}
\definecolor{codegreen}{rgb}{0,0.6,0}
\definecolor{codegray}{rgb}{0.5,0.5,0.5}
\definecolor{codepurple}{rgb}{0.58,0,0.82}
\definecolor{backcolour}{rgb}{0.95,0.95,0.92}
\lstdefinestyle{cstyles}{
language=C,
backgroundcolor=\color{backcolour},
commentstyle=\color{codegreen},
keywordstyle=\color{magenta},
numberstyle=\tiny\color{codegray},
stringstyle=\color{codepurple},
basicstyle=\ttfamily\footnotesize,
breakatwhitespace=false,
breaklines=true,
captionpos=b,
keepspaces=true,
numbers=left,
numbersep=5pt,
showspaces=false,
showstringspaces=false,
showtabs=false,
tabsize=2,
extendedchars=true,
literate=
  {&&}{{\&\&}}2
{||}{{||}}2
{!=}{{$\neq$}}2
{==}{{==}}2
{->}{$\rightarrow$}2
{<=}{{$\leq$}}2
{>=}{{$\geq$}}2
{á}{{\'a}}1
{ã}{{\~a}}1
{à}{{\`a}}1
{â}{{\^a}}1
{é}{{\'e}}1
{ê}{{\^e}}1
{í}{{\'i}}1
{ó}{{\'o}}1
{õ}{{\~o}}1
{ô}{{\^o}}1
{ú}{{\'u}}1
{*}{{$^*$}}1
{ç}{{\c{c}}}1
}

\title{Aplicação de HASH e Filtro de Bloom}
\subtitle{Busca de Números Primos}

\institute{UFERSA}
\author{Prof. Kennedy Lopes}
\date{\today}

\newcounter{meucontador} % Cria o contador "meucontador"
\newcommand{\novoitem}{\refstepcounter{meucontador}\arabic{meucontador}}

\begin{document}

\begin{frame}
  \titlepage
\end{frame}

\newcommand{\addsection}[2]{\section{#1}\input{#2}}
\newcommand{\addsubsection}[2]{\subsection{#1}\input{#2}}

\begin{frame}
  \frametitle{Proposta}
  Um sistema embarcado possui pouca memória e precisa responder rapidamente consultas sobre números primos.
  Objetivos:
  \begin{itemize}
    \item Gerar os números primos até \texttt{$MAX\_PRIMO=1000$};
    \item Inserir os primos no filtro de Bloom;
    \item Consutar inteiros;
    \item Informar sobre o resultado probabilístico.
  \end{itemize}
\end{frame}

\begin{frame}
  \frametitle{Crivo de Erastóstenes}

  O método consiste em listar todos os números naturais a partir do 2 até o valor limite desejado e seguir um passo a passo progressivo:
  \begin{enumerate}
    \item O primeiro número primo é o 2. Ele é o primeiro número primo. Eliminar múltiplos do 2
    \item Mantém-se o 2, mas riscam-se todos os seus múltiplos (4, 6, 8, 10, etc.).
    \item O próximo número não riscado após o 2 é o 3, que também é primo.Eliminar múltiplos do 3;
    \item Mantém-se o 3 e riscam-se todos os seus múltiplos restantes (9, 15, 21, etc.).
    \item Repete-se o processo para o próximo número não riscado (o 5, depois o 7, e assim por diante).
    \item Quando se atinge um número cujo quadrado é maior que o limite da tabela, todos os números restantes que não foram riscados são primos.
  \end{enumerate}

\end{frame}

\begin{frame}
  \frametitle{Geração dos primos}
  \begin{multicols}{2}
    \lstinputlisting[style=cstyles, numbers=none, frame=shadowbox, basicstyle=\ttfamily\small]{../codes/geracaoPrimos.c}
  \end{multicols}
\end{frame}



\begin{frame}
  \frametitle{Estrutura do Filtro}

  \begin{itemize}
    \item \texttt{struct} contendo um vetor de bytes;
    \item Cada byte pode guardar até 8 valores alvos de \emph{hash}.
    \item Necessário indicar onde está cada valor;
  \end{itemize}

\end{frame}
\begin{frame}
  \frametitle{Implementação}
  \lstinputlisting[style=cstyles, numbers=none, frame=shadowbox, basicstyle=\ttfamily\tiny]{../codes/main.c}
\end{frame}

\begin{frame}
  \frametitle{Estrutura de Bytes}

  \lstinputlisting[style=cstyles, numbers=none, frame=shadowbox, basicstyle=\ttfamily\tiny]{../codes/main.c}

\end{frame}

\begin{frame}
  \frametitle{Mapeando byte}

\end{frame}

\begin{frame}
  \frametitle{Mapeando bit}

\end{frame}

\begin{frame}
  \frametitle{Funções HASH}


\end{frame}

\begin{frame}
  \frametitle{Ponteiros de Funções}



\end{frame}

\begin{frame}
  \frametitle{Analisando filtro}

\end{frame}


\begin{frame}
  \frametitle{Checando número primo}

\end{frame}

\begin{frame}
  \frametitle{Falso Positivos x Falsos Negativos}

\end{frame}

\begin{frame}[allowframebreaks]
  \frametitle{Otimizações}
  Mater informações atualizadas como \emph{Falsos Positivos} e \emph{Taxa de ocupação do filtro}.
  Parâmetros:
  \begin{itemize}
    \item \(n\): Número de bits
    \item \(k\): Número de funções Hash
    \item \(m\): Elementos Inseridos
  \end{itemize}

  Probabilidade de um bit permanecer \(0\) após n inserções:

  \[P_0 = \left(1-\frac{1}{m}\right)^{kn} \approx e^{-kn/m}\]

  Probabilidade de ser 1 após n inserções:

  \[P_1 = 1 - P_0\]

  Probabilidade de Falso positivos:

  \[P_{fp} = P_1^k\]

  Otimizando \(k\) mantendo \(m\) e \(n\) fixos:

  $$k_{otimo} = \frac{m}{n} \ln(2)$$

  Conhecendo o $k_{otimo}$, podemos dimensionar o filtro para os valores de \(m\) ou \(n\).
\end{frame}

\begin{frame}
  \frametitle{Propondo estruturas}

\end{frame}

\end{document}
